\sezione{6}{La normale univariata e la stima MLE}{26–30}

Inizia l'approccio parametrico. Si parte dal caso più semplice, una sola feature ($d=1$), per capire \emph{com'è fatta} la gaussiana e \emph{come si stimano} i suoi parametri dai dati.

\subsection{La distribuzione normale}
\begin{formulat}{Densità normale ($d=1$)}
$\displaystyle p(x)=\frac{1}{\sigma\sqrt{2\pi}}\;\exp\!\left[-\frac{(x-\mu)^2}{2\sigma^2}\right]$
\end{formulat}
\begin{intuizione}
È la classica «campana»: la maggior parte dei valori si concentra attorno a un centro e quelli estremi diventano sempre più rari. Bastano \textbf{due numeri}: $\mu$ (media) dice \emph{dove} sta la campana, $\sigma$ (deviazione standard) quanto è \emph{larga}; $\sigma^2$ è la varianza.
\end{intuizione}

\begin{center}
\rowcolors{2}{chiaro}{white}
\begin{tabularx}{\linewidth}{>{\centering\arraybackslash}m{3cm} Y}
\rowcolor{navy}\thd{Pezzo} & \thd{Significato}\\
$(x-\mu)$ & distanza di $x$ dal centro\\
$(x-\mu)^2/\sigma^2$ & distanza al quadrato \textbf{misurata in unità di $\sigma$}: essere a 10 dal centro è «tanto» se $\sigma=2$, «poco» se $\sigma=50$. Diventerà la distanza di Mahalanobis\\
$\exp(-\,\cdot/2)$ & trasforma la distanza in plausibilità: distanza 0 → massimo; distanza grande → quasi 0\\
$1/(\sigma\sqrt{2\pi})$ & normalizzazione: area totale $=1$. Se $\sigma$ cresce, la campana si allarga e deve \emph{abbassarsi}\\
\end{tabularx}
\end{center}

\fig[.8\linewidth]{gauss_regola.pdf}
\didascalia{Regola 68–95–99.7: oltre $3\sigma$ il contributo è trascurabile (la gaussiana però non tocca mai lo zero).}

\begin{attenzione}[Il N.B. della slide 26]
Questa $p(x)$ è la densità di \emph{una} gaussiana; nel classificatore diventerà una $p(x|w_i)$, la densità di una classe. Non va confusa con la $p(\x)$ al denominatore di Bayes, che è la \emph{miscela} $\sum_i p(\x|w_i)P(w_i)$ di tutte le classi.
\end{attenzione}

\subsection{L'idea della Massima Verosimiglianza (MLE)}
\begin{intuizione}
Abbiamo 10 numeri: $\{3,\,7,\,9,\,-2,\,15,\,54,\,-11,\,0,\,23,\,-8\}$ e assumiamo che vengano da una gaussiana. \emph{Quale}? Una campana centrata su 100 rende i dati quasi impossibili; una strettissima centrata su 9 spiega il 9 ma non il 54. C'è una campana che rende i dati osservati «\textbf{più probabili possibile}»: è la stima MLE.
\end{intuizione}
\begin{formulat}{Funzione di verosimiglianza (campioni indipendenti)}
$\displaystyle L(\mu,\sigma)=\prod_{i=1}^{n} p(x_i\,|\,\mu,\sigma)
\qquad\qquad (\hat\mu,\hat\sigma)=\argmax_{\mu,\sigma} L(\mu,\sigma)$
\end{formulat}
Cambio di prospettiva: prima $p(x|\mu,\sigma)$ fissava i parametri e guardava $x$; ora i \textbf{dati sono fissi} (osservati) e variano i \textbf{parametri}.

\begin{mlbox}[MLE e MAP: due argmax diversi]
MAP (sezione 2): $\argmax$ sulle \emph{classi} di $P(\text{classe}|\text{dati})$, usa il prior. MLE: $\argmax$ sui \emph{parametri} di $p(\text{dati}|\text{parametri})$, senza prior. Aggiungendo un prior sui parametri si ottiene la stima MAP dei parametri.
\end{mlbox}

\subsection{Da dove vengono le formule}
\begin{tecnica}[Derivazione in tre passi]
\textbf{1. Logaritmo.} Massimizzare $L$ equivale a massimizzare $\ln L$ (il log è crescente) e il prodotto diventa una somma; in pratica evita anche l'\emph{underflow} numerico:
\[
\ell(\mu,\sigma^2)=-\frac n2\ln(2\pi\sigma^2)-\frac1{2\sigma^2}\sum_{i=1}^{n}(x_i-\mu)^2
\]
\textbf{2.} $\dfrac{\partial\ell}{\partial\mu}=0\;\Rightarrow\;\sum_i(x_i-\mu)=0\;\Rightarrow\;\hat\mu=\dfrac1n\sum_i x_i$\qquad (media campionaria)\\[4pt]
\textbf{3.} $\dfrac{\partial\ell}{\partial\sigma^2}=0\;\Rightarrow\;\hat\sigma^2=\dfrac1n\sum_i (x_i-\hat\mu)^2$\qquad (varianza campionaria)
\end{tecnica}
Intuizione del passo 2: massimizzare la likelihood gaussiana equivale a \textbf{minimizzare la somma dei quadrati} delle distanze dal centro, come nei minimi quadrati.

\subsection{L'esempio numerico}
\begin{esempio}[Stima di $\mu$ e $\sigma$ con $n=10$]
\[
\hat\mu=\frac{3+7+9-2+15+54-11+0+23-8}{10}=\frac{90}{10}=9
\]
\begin{center}\small
\begin{tabular}{l*{10}{c}}
\toprule
$x_i$ & 3 & 7 & 9 & $-2$ & 15 & \textcolor{rosso}{\textbf{54}} & $-11$ & 0 & 23 & $-8$\\
$x_i-9$ & $-6$ & $-2$ & 0 & $-11$ & 6 & \textcolor{rosso}{\textbf{45}} & $-20$ & $-9$ & 14 & $-17$\\
$(x_i-9)^2$ & 36 & 4 & 0 & 121 & 36 & \textcolor{rosso}{\textbf{2025}} & 400 & 81 & 196 & 289\\
\bottomrule
\end{tabular}
\end{center}
\[
\hat\sigma^2=\frac{3188}{10}=318.8\qquad\qquad \hat\sigma=\sqrt{318.8}\approx17.855
\]
(Nella slide 28 c'è un «$^2$» in più dopo l'ultima parentesi del numeratore: è un refuso, il risultato è corretto.)
\end{esempio}

\fig[.98\linewidth]{mle1d.pdf}
\didascalia{Sinistra: la gaussiana stimata e la densità nel punto 25. Destra: la log-verosimiglianza in funzione di $\mu$ ha il massimo proprio nella media campionaria.}

\begin{attenzione}[Due osservazioni importanti]
\textbf{a) L'outlier domina.} Il valore 54 contribuisce 2025 su 3188, cioè il \textbf{64\%} della varianza. Togliendolo: $\hat\mu=4$, $\hat\sigma\approx10.2$. Il MLE gaussiano è sensibile agli outlier perché eleva le distanze al quadrato.

\textbf{b) $n$ o $n-1$?} La MLE divide per $n$ ed è leggermente \emph{distorta} (sottostima). Dividendo per $n-1$ (correzione di Bessel) si ha la stima corretta: $3188/9\approx354.2$. La differenza conta solo con pochi dati. Il prof usa $1/n$, coerente con la MLE.
\end{attenzione}

\subsection{Densità $\neq$ probabilità}
\begin{formula}
$\displaystyle p(25)=\frac{1}{17.855\sqrt{2\pi}}\;e^{-\frac{(25-9)^2}{2\cdot318.8}}=\frac{1}{44.7559}\,e^{-0.4015}\approx0.01495$
\end{formula}
\begin{intuizione}
$p(25)$ \textbf{non} è «la probabilità di osservare esattamente 25»: nel continuo quella è 0. È una \emph{concentrazione} di probabilità attorno a 25. Per avere una probabilità si moltiplica per la larghezza di un intervallo:
\[
P(x\le X\le x+dx)\approx p(x)\,dx
\]
Come la densità di un materiale (kg/m$^3$) non è una massa finché non la si moltiplica per un volume.
\end{intuizione}
\begin{esempio}[Verifica numerica]
La probabilità esatta dell'intervallo $[24.5,\,25.5]$ (largo 1) vale $0.014954$, praticamente uguale a $p(25)\cdot1=0.014955$. Controllo di sanità: in $\hat\mu\pm2\hat\sigma=[-26.7,\,44.7]$ cadono 9 valori su 10; resta fuori solo il 54.
\end{esempio}
Una densità \textbf{può superare 1}: con $\sigma=0.1$ il picco vale circa $3.99$. Conta solo che l'area totale sia 1. Nella formula di Bayes il $dx$ si semplifica e la posteriori resta una vera probabilità.

\begin{appunti}{— Sezione 6}
\pag{26} $\mu$ = posizione, $\sigma$ = larghezza; $(x-\mu)^2/\sigma^2$ = distanza in unità di $\sigma$. $1/(\sigma\sqrt{2\pi})$: area 1. 68/95/99.7\% in $\pm1,2,3\sigma$. Questa $p(x)$ diventerà $p(x|w_i)$, $\neq$ evidence.
\pag{27} MLE: parametri che rendono i dati più verosimili. $L=\prod_i p(x_i|\mu,\sigma)$; dati fissi, parametri variabili; si massimizza $\ln L$. MLE (parametri, nessun prior) vs MAP (classi, con prior).
\pag{28} $\hat\mu=\bar x=9$; $\hat\sigma^2=318.8$; $\hat\sigma\approx17.855$. Outlier 54 = 64\% della varianza. MLE divide per $n$ (distorta); $n-1$: 354.2. Refuso: «$^2$» finale.
\pag{29–30} $p(25)\approx0.01495$ è una densità; $P([x,x+dx])\approx p(x)dx$; densità $>1$ possibile; in Bayes il $dx$ si semplifica.
\end{appunti}
